\title{Can Intelligence Be Wise?}
\begin{document}
\maketitle

\begin{abstract}
Can Intelligence Be Wise? Hakan Altinay 28 May 2026 Health and Social Policy Hakan Altinay on why listening and love may be more important than intelligence. First century thinker Seneca is supposed to have noted that if one does not know to which port one should sail, no wind is favorable. Bill Gates, for his part, tells us that innovation accelerated by AI will mitigate many of our contemporary problems, including the effects of climate change. Others highlight stupidity as the main problem we suffer from as humanity and artificial general intelligence as the solution towards which we should all strive and rejoice. That particular diagnosis about stupidity by Yann LeCun is certainly in line with the general Silicon Valley ethos which prides itself for channeling abundant IQ to tasks at hand. So, we know their direction of sail and favored winds. There has been more than one reason to be suspicious of the diagnosis and the proposed treatment protocol above. Were these not the people who told us to move fast and break things? We are now reassured that AI will only diffuse at the speed of trust. Only a decade ago, social networks were sold as the seamless way to be closer to friends, family, and friends-to-be, only to discover that these platforms allow us to see a meager 7\% of content from our chosen connections. These fake paragons of connections became obscenely rich along the way. Fool me once, shame on you; fool me twice, shame on me. My own unease lies at a somewhat diff
\end{abstract}

\section{Recovered Text}
Can Intelligence Be Wise?
Hakan Altinay
28 May 2026
Health and Social Policy
Hakan Altinay on why listening and love may be more important than intelligence.
First century thinker Seneca is supposed to have noted that if one does not know to which port one should
sail, no wind is favorable. Bill Gates, for his part, tells us that innovation accelerated by AI will mitigate
many of our contemporary problems, including the effects of climate change. Others highlight stupidity as
the main problem we suffer from as humanity and artificial general intelligence as the solution towards which
we should all strive and rejoice. That particular diagnosis about stupidity by Yann LeCun is certainly in line
with the general Silicon Valley ethos which prides itself for channeling abundant IQ to tasks at hand. So, we
know their direction of sail and favored winds.
There has been more than one reason to be suspicious of the diagnosis and the proposed treatment protocol
above. Were these not the people who told us to move fast and break things? We are now reassured that AI
will only diffuse at the speed of trust. Only a decade ago, social networks were sold as the seamless way to be
closer to friends, family, and friends-to-be, only to discover that these platforms allow us to see a meager 7\%
of content from our chosen connections. These fake paragons of connections became obscenely rich along
the way. Fool me once, shame on you; fool me twice, shame on me.
My own unease lies at a somewhat different coordinate. For years, I worked for a large international foun-
dation, and one of the wisest people I knew described a senior member of that foundation as being “too
smart for him.” Since I have not sought and received permission to disclose the names in question, I should
simply report that cumulative actions of the wise friend helped transform a very troublesome cross border
relationship, while the work of the foundation and that colleague amounted to rather little in the end. So, the
question becomes whether intelligence is really intelligent or wise.
Sezai Karakoç, a Turkish poet, argues that love is wiser than intelligence. Others have agreed with him: John
O’Donohue, this time an Irish poet, notes that love is the only light which allows us to really see the other.
In the latest encyclical from the Vatican, we are reminded that “when we grasp a reality with our heart, we
know it better and more fully.” The same encyclical warns us that devaluing the heart would mean knowing
too narrowly and missing out on, among others, poetry. James Baldwin, who cannot be accused of being
easy prey, observed that the longer he lived, the more deeply he learnt that love as friendship was the work
of realizing each other’s light.
Simon Critchley, an irreverent contemporary philosopher, has a related challenge which should concern many
of us: If overthinking means underloving, do we still want to overthink?
Love, in this account, is of course not the love narrated in Hollywood movies but what Iris Murdoch would
describe as the quality of the attention we pay. Listening well, for example, has been called an act of love.
Several traditions, such as the Benedictines, Taoists and Sufis, have the phrase “listen with your heart’s ear”.
1
That phrasing may be predicated on the awareness that to really hear the other, you need to intend to hear the
other and not merely wait for your turn to speak.
What seems to be at stake is not yet another gimmick to optimize and self-improve, but an existential choice:
Is the other a potential foe or competitor? Or, is s/he another you whom you have not yet befriended? Are
they a means or an end in themselves?
Listening with your heart’s ear is, along with curiosity, the central rule we agree upon at the outset of our
programs at the European School of Politics. Those programs attract very diverse participants and lead to
deeper and stronger relationships than anybody expects. The fact that such a modest effort leads to such
profound changes should theoretically interest Yann LeCun and his peers, but we cannot bring ourselves to
expect this from them.
Iain McGilchrist is one helpful scholar as we try to come to terms with the seeming shallowness of what gen-
erally passes as intelligence and effectiveness of its opposite. McGilchrist describes two modes of knowing,
one that is narrow and instrumental and another which is holistic and hermeneutic. He describes the way we
move through the world as the interplay of these two different modes and hemispheres. One mode catego-
rizes in a rush and manipulates, the other attends and relates. Our contemporary challenge is that we have
overdeveloped the former at the expense of the latter. Love has occasionally been called a muscle. Could it
be that that vital muscle has atrophied? And if so, what would the intelligence that will emanate on steroids
from the current constellation end up serving? Will it make us wiser, or more feral and therefore less capable
of wisdom that Karakoç and company refer to?
Fortunately, every action causes a reaction. There are everywhere people whom Maria Popova describes
with her usual and profound insights: those “who alchemize suffering and rage into love, who compost
disappointment into fertilizer for growth, who break down cynicism to its building blocks of helplessness
and hubris, then metabolize the toxin out of the system we call society.” I met one of those people recently
at the Frontiers of Democracy Conference at Tufts University: Ken Powley. Mr. Powley has spent 50 years
organizing whitewater rafting expeditions, only to change course after the January 6 raid on the Capitol Hill
and establish Team Democracy, which uses Mr. Powley’s earlier networks in his industry to provide rafting
experience to people from divergent political tribes. His conviction is that being literally on the same boat
under testing conditions provides a reality check and a somatic realization of the stakes.
Fortunately, Mr. Powley is not alone: Big Lunch has been going on in the UK for 15 years. Weavers in the
US, whose motto includes love, has a 10+ year history. One Radical Centre Manifesto invites all to “a place
of encounter that swims against the tide of polarisation. It is radical not because it is extreme, but because it
resists the forces pulling us apart. And it is a centre not in the sense of a political midpoint or compromise
platform, but as a shared, imagined ground where we agree to meet one another—across difference, uncer-
tainty, and disagreement—for the sake of deeper, shared values.” The good colleagues at the St. Ethelburga’s
Centre note that the Radical Centre is not a political program, a set of policies, or a new ideological brand,
or a hidden attempt to smuggle one worldview past another, but a practice and a muscle. Harsh Mander,
an Indian sage and a possible heir to the Gandhian practice and muscle, describes what he does as radical
love and is determined to leave no ostracized Muslim unheeded and unaccompanied. Kumi Naidoo, also no
stranger to struggle, argues that if we do not find a way to love Trump voters that would be our insufficiency
and not theirs. I cannot help wondering whether Yann LeCun’s enhanced intelligence has—or may have in
the future—a way to account for such dispositions or would simply treat them as maladaptive choices or
mistakes in the way many economists have called generous choices in the ultimatum game experiments.
To recap, humanity’s main shortcoming may not be stupidity and greater instrumental intelligence may not
be the solution. There might be a bigger question we have to answer: Are we, when all is said and done,
atomistic units with negligible ties to our natural and social habitats, which are there primarily for use and
2
extraction? Do we have any other north star other than more consumption and more power? Or, are we
constitutively—and fortunately—interdependent in which case care, love, goodwill are not trivial but central
to our wellbeing and flourishing? My sense is that we have neglected our wherewithal to generate and sustain
temperaments such as care, goodwill and love for a very long time. A civilization with meager capacities
to really see and cherish its members may end up becoming less competent and wise by becoming more
narrowly intelligent and powerful.
None of this means that the corrective I propose is easy, self-evident, or assured of success. I myself was
a product of a Bildung which cherished critical thinking as the obvious standard, and yet I am now thrilled
that my son attends a school whose motto is “with love we flourish.” My hunch is that people such as
Kumi Naidoo, Harsh Mander, the wise friend I mentioned at the outset, many ordinary yet kind people, and
of course Hrant Dink were the invaluable breadcrumbs which helped me discover somatically a different
paradigm. This may have been one instance of the friendship as a theory of change phenomenon. I may
have opted for a different track because I couldn’t bring myself to unknow the invaluable realizations I was
fortunate enough to know through these people. And if that is the case, an ecosystem which can produce,
sustain and enhance the sort of people that Maria Popova described, becomes ever more important.
Once again, we are both the problem and the elusive solution at the same time. Time will tell whether enough
of us choose to walk that path with no guarantees of success and ample evidence for probable failure.
Hakan Altinay is the founding Director of the European School of Politics in Istanbul, and a Professor of the
Practice at Tufts University in Boston.
3
\end{document}
